\chapter{The Deligne--Laumon Direct-Image Theorem}
\label{chap:global-weights}

We now pass from the estimate on a single curve to the global theorem on
weights.  The reduction has two independent parts.  Geometry and
Noetherian induction reduce an arbitrary morphism to a relative affine
curve.  Proper-support base change then identifies the stalks of its
higher direct images with the compactly supported cohomology treated in
\cref{thm:curve-estimate}.  Mixedness is not inferred merely from
pointwise inequalities: it is obtained from real envelopes and
\cref{thm:real-sheaf-mixed}.  This separation is essential.

\section{The theorem}

\begin{theorem}[Deligne--Laumon direct-image theorem]
\label{thm:deligne-laumon-direct-image}
Let
\[
  f_0:X_0\longrightarrow Y_0
\]
be a separated morphism of finite type between schemes of finite type
over \(k_0\).  If a constructible Weil sheaf \(\cF_0\) on \(X_0\) is
\(\iota\)-mixed of weights at most \(w\), then
\[
  R^i f_{0!}\cF_0
\]
is \(\iota\)-mixed of weights at most \(w+i\) for every \(i\).

More generally, if a constructible Weil complex \(\cK_0\) has weights at
most \(w\), in the sense that
\[
  \cH^j(\cK_0)\text{ has weights at most }w+j,
\]
then \(Rf_{0!}\cK_0\) has weights at most \(w\).
\end{theorem}

The theorem is local on \(Y_0\), but it is not enough to inspect a single
fiber of an arbitrary-dimensional morphism: that would ask for the very
theorem we are proving.  We first reduce the relative dimension to one.

\section{Formal reductions}

We list the reductions as lemmas so that none is hidden in the phrase
``standard devissage''.

\begin{lemma}[Quasi-finite case]
\label{lem:weights-quasifinite}
If \(f_0\) is separated and quasi-finite, then \(R^if_{0!}=0\) for
\(i>0\), and \(f_{0!}\) preserves upper weight bounds.
\end{lemma}

\begin{proof}
By Zariski's main theorem, locally on the target \(f_0\) factors as an
open immersion followed by a finite morphism.  Extension by zero is
exact, and finite pushforward is exact, so the higher direct images
vanish.  At a geometric point of the target, \(f_!\cF\) is the direct
sum of stalks over the finite fiber, with zero contributions from points
outside the chosen open.  Frobenius permutes the summands in finite
orbits.  The characteristic polynomial of one orbit is obtained from
the Frobenius polynomial at the corresponding closed point upstairs by
replacing \(T\) by a power.  As in the finite-map reduction in
\cref{chap:curve-estimate}, the normalized absolute value and the residue
degree are multiplied by the same integer.  Thus weights do not
increase.
\end{proof}

\begin{lemma}[Composition]
\label{lem:weights-composition}
Suppose the conclusion of
\cref{thm:deligne-laumon-direct-image} is known for
\(h_0:X_0\to Z_0\) and \(g_0:Z_0\to Y_0\).  Then it is known for
\(f_0=g_0h_0\).
\end{lemma}

\begin{proof}
The compact-support Leray spectral sequence is Frobenius-equivariant:
\[
  E_2^{a,b}
  =
  R^ag_{0!}R^bh_{0!}\cF_0
  \Longrightarrow
  R^{a+b}f_{0!}\cF_0.
\]
If \(\cF_0\) has weights at most \(w\), the first application gives
weights at most \(w+b\), and the second gives weights at most
\(w+b+a\).  The abutment has a finite Frobenius-stable filtration whose
graded pieces are subquotients of \(E_2^{a,b}\).  Weight bounds pass to
subquotients and extensions by
\cref{prop:elementary-weight-permanence}.  No degeneration of the
spectral sequence is asserted or needed.
\end{proof}

\begin{lemma}[Open--closed devissage]
\label{lem:weights-open-closed}
Let \(j_0:U_0\hookrightarrow X_0\) be open with closed complement
\(i_0:Z_0\hookrightarrow X_0\).  If the theorem is known for the
restrictions of \(f_0\) to \(U_0\) and \(Z_0\), then it is known for
\(f_0\).
\end{lemma}

\begin{proof}
Use the distinguished triangle
\[
  j_{0!}j_0^*\cK_0
  \longrightarrow\cK_0
  \longrightarrow i_{0*}i_0^*\cK_0
  \xrightarrow{+1}.
\]
After applying \(Rf_{0!}\), the long exact cohomology sequence expresses
each cohomology sheaf as an extension of a quotient and a subobject of
the two restricted terms.  Upper weight bounds survive all three
operations.
\end{proof}

\begin{lemma}[Generic relative Noether normalization]
\label{lem:relative-noether-normalization}
After replacing \(X_0\) and \(Y_0\) by dense opens on each irreducible
component, a dominant affine morphism \(f_0:X_0\to Y_0\) of relative
dimension \(r\) admits a factorization
\[
 X_0\xrightarrow{\nu_0}\A^r_{Y_0}
 \xrightarrow{\pi_0}Y_0
\]
with \(\nu_0\) finite.  The projection \(\pi_0\) is a composition of
\(r\) relative affine-line projections.
\end{lemma}

\begin{proof}
Apply Noether normalization to the finitely generated extension of the
function field of \(Y_0\) by that of \(X_0\).  The resulting algebraically
independent functions and the integral dependence relations spread out
after inverting finitely many functions on \(X_0\) and \(Y_0\).  They
define the finite map to affine \(r\)-space.  The last assertion is the
coordinate-by-coordinate factorization of affine space.
\end{proof}

\section{Reduction to a relative affine curve}

We prove the sheaf statement of
\cref{thm:deligne-laumon-direct-image} by simultaneous Noetherian
induction on the supports in source and target.

Use a finite stratification and a weight filtration to reduce to
\(\cF_0\) pure and lisse on a smooth, irreducible, locally closed
stratum of \(X_0\).  Extension by zero and
\cref{lem:weights-open-closed} allow us to replace \(X_0\) by that
stratum.  Shrinking the target to the dense image, and treating the
closed complement by induction, we may assume that \(f_0\) is dominant
and both source and target are affine.  Relative Noether normalization
and \cref{lem:weights-quasifinite,lem:weights-composition} reduce the
problem to a projection
\[
  \pi_0:\A^1_{Y_0}\longrightarrow Y_0.
\]

The finite pushforward from the normalization step is constructible, not
necessarily lisse everywhere.  Stratify \(\A^1_{Y_0}\), use
\cref{lem:weights-open-closed}, and shrink the base once more.  We arrive
at
\[
  h_0:U_0\longrightarrow Y_0,
\]
where \(U_0\) is a dense open of \(\A^1_{Y_0}\), \(h_0\) is smooth and
surjective, all geometric fibers are nonempty geometrically irreducible
affine curves, and \(\cF_0\) is lisse and pure of weight \(w\).

All discarded closed subsets have strictly smaller support and are
covered by the induction.  This proves that it is enough to establish
the theorem in the relative affine-curve case.

\section{Fiberwise upper bounds}

Let \(y_0\) be a closed point of \(Y_0\), with residue field
\(k(y_0)\) of cardinality \(q_y=q^{\deg(y_0)}\), and choose a geometric
point \(\bar y\) above it.  Proper-support base change gives a
Frobenius-equivariant isomorphism
\[
  \left(R^ih_{0!}\cF_0\right)_{\bar y}
  \simeq
  H_c^i(U_{\bar y},\cF|_{U_{\bar y}}).
\tag{12.1}
\label{eq:relative-curve-base-change}
\]
The restriction of \(\cF_0\) to the fiber is pure of weight \(w\)
relative to \(q_y\): a degree \(n\) point over \(k(y_0)\) is a degree
\(n\deg(y_0)\) point over \(k_0\).  Applying
\cref{thm:curve-estimate} over \(k(y_0)\) shows that every eigenvalue in
\eqref{eq:relative-curve-base-change} has weight at most \(w+i\).
Thus \(R^ih_{0!}\cF_0\) is punctually of weights at most \(w+i\).

Punctual bounds do not by themselves construct a mixed filtration on
the base.  We now supply it without invoking the theorem recursively.

\section{Mixedness in the relative-curve case}

\begin{lemma}
\label{lem:relative-curve-mixedness}
In the relative affine-curve situation above,
\[
  R^ih_{0!}\cF_0
\]
is \(\iota\)-mixed of weights at most \(w+i\).
\end{lemma}

\begin{proof}
It is enough first to treat a pure \(\cF_0\).  Form its real envelope
\[
  \cE_0=\cF_0\oplus\cF_0^c
\]
as in \cref{lem:real-envelope}; both summands are pure of weight \(w\),
and \(\cE_0\) is \(\iota\)-real.  Since \(R^ih_!\cF\) is a direct
summand of \(R^ih_!\cE\), it will be enough to prove that the latter is
mixed.

For a geometric affine-curve fiber, \(H_c^0=0\).  Top cohomology is
\[
 H_c^2(U_{\bar y},\cE)
 \simeq
 \left(E_{\bar\eta}\right)_{\pi_1(U_{\bar y})}(-1).
\]
The coinvariant quotient is geometrically constant.  Since \(\cE_0\)
is pure of weight \(w\), every nonzero eigenvalue on this top cohomology
has weight exactly \(w+2\).  By
\cref{thm:curve-estimate}, the eigenvalues on \(H_c^1\) have weights at
most \(w+1\).

The fiber \(L\)-function of the real sheaf \(\cE_0\) has real
coefficients after \(\iota\), because every local factor has real
coefficients.  The trace formula writes it as
\[
  L(U_{y_0},\cE_0,T)
  =
  \frac{
    \det(1-TF_y\mid H_c^1)
  }{
    \det(1-TF_y\mid H_c^2)
  }.
\tag{12.2}
\label{eq:relative-real-fiber-L}
\]
Roots of the numerator have reciprocal absolute value at least
\(q_y^{-(w+1)/2}\), whereas every root of the denominator has reciprocal
absolute value exactly \(q_y^{-(w+2)/2}\).  The two polynomials are
coprime.  Since their quotient is fixed by complex conjugation and both
have constant term \(1\), uniqueness of a reduced fraction shows that
each polynomial has real coefficients.  This remains true after every
finite extension of \(k(y_0)\), so the local characteristic polynomials
of \(R^1h_!\cE\) and \(R^2h_!\cE\) are \(\iota\)-real.

After stratifying \(Y_0\), the higher direct images are lisse on each
stratum.  Apply \cref{thm:real-sheaf-mixed} on each stratum and use
open--closed induction to conclude that they are mixed.  Their punctual
upper bounds were already proved from \cref{eq:relative-curve-base-change},
so their mixed weights are at most \(w+i\).  Direct summands inherit
mixedness and the bound, giving the assertion for \(\cF_0\).

For a general mixed \(\cF_0\), choose a finite weight filtration.  Apply
\(R^ih_!\) to each short exact sequence.  The resulting long exact
sequences, together with induction on the filtration length and
subquotient/extension permanence, prove the claim.
\end{proof}

\begin{warning}
The coprimality in \eqref{eq:relative-real-fiber-L} uses the strict gap
between \(w+1\) and \(w+2\).  Merely knowing that the whole alternating
product is real would not imply that each cohomological factor is real.
This is the only no-cancellation argument needed to pass from punctual
bounds to mixed sheaves in the reduction.
\end{warning}

\section{Completion of the direct-image theorem}

\begin{proof}[Proof of \cref{thm:deligne-laumon-direct-image}]
For a sheaf, the formal reductions of the preceding sections reduce the
claim to the relative affine-curve case, which is
\cref{lem:relative-curve-mixedness}.  Open--closed induction restores the
discarded strata, composition restores the successive affine-line
projections, and the quasi-finite case restores the finite normalization
map.  This proves the sheaf statement.

Let now \(\cK_0\) be a bounded constructible complex of weights at most
\(w\).  The hypercohomology spectral sequence
\[
  E_2^{a,b}
  =
  R^af_{0!}\cH^b(\cK_0)
  \Longrightarrow
  \cH^{a+b}(Rf_{0!}\cK_0)
\tag{12.3}
\label{eq:weight-hypercohomology-ss}
\]
is Frobenius-equivariant and finite.  The sheaf
\(\cH^b(\cK_0)\) has weights at most \(w+b\), so the result just proved
puts weights at most \(w+a+b\) on \(E_2^{a,b}\).  The abutment in total
degree \(n=a+b\) has a finite filtration by subquotients of those terms
and therefore has weights at most \(w+n\).  This is precisely the derived
weight convention from \cref{def:iota-purity}.
\end{proof}

\section{Smooth proper purity}

Upper bounds become exact under smoothness and properness because
relative Poincare duality supplies the opposite bound.

\begin{corollary}[Smooth-proper purity]
\label{cor:smooth-proper-purity}
Let \(f_0:X_0\to Y_0\) be smooth and proper of pure relative dimension
\(d\), and let \(\cF_0\) be a lisse Weil sheaf punctually
\(\iota\)-pure of weight \(w\).  Then \(R^if_{0*}\cF_0\) is lisse and
punctually \(\iota\)-pure of weight \(w+i\) for every \(i\).
\end{corollary}

\begin{proof}
Properness identifies \(Rf_!=Rf_*\), and smooth proper base change makes
each \(R^if_*\cF\) lisse.  The direct-image theorem gives weights at most
\(w+i\).

Relative Poincare duality gives a Frobenius-equivariant isomorphism
\[
  (R^if_*\cF)^\vee
  \simeq
  R^{2d-i}f_*(\cF^\vee)(d).
\tag{12.4}
\label{eq:relative-pd-weight}
\]
The dual sheaf \(\cF^\vee\) is pure of weight \(-w\).  Applying the upper
bound in degree \(2d-i\) and then twisting by \((d)\) gives weights at
most
\[
  -w+(2d-i)-2d=-w-i
\]
on the right side of \eqref{eq:relative-pd-weight}.  Dualizing reverses
weights, so \(R^if_*\cF\) has weights at least \(w+i\).  Combining upper
and lower bounds proves purity.
\end{proof}

\begin{corollary}[Cohomological Riemann hypothesis]
\label{cor:cohomological-rh}
Let \(X_0\) be smooth and proper over \(k_0\).  For every \(i\), every
eigenvalue \(\alpha\) of geometric Frobenius on
\[
  H^i(X,\Qbarl)
\]
satisfies
\[
  |\iota(\alpha)|=q^{i/2}.
\]
The assertion holds for every isomorphism
\(\iota:\Qbarl\to\mathbf C\).
\end{corollary}

\begin{proof}
Apply \cref{cor:smooth-proper-purity} to
\(X_0\to\Spec k_0\) and the constant sheaf, which is pure of weight
\(0\).  The entire argument from \cref{chap:weights} onward was valid
for an arbitrary choice of \(\iota\), so the conclusion may be repeated
for every such choice.
\end{proof}

\begin{remark}
At this stage the Frobenius eigenvalues are elements of \(\Qbarl\) with
the asserted absolute value under every field isomorphism to
\(\mathbf C\).  The next part of the notes combines this with the
rational zeta function to prove that they are algebraic over
\(\mathbf Q\), that every complex conjugate has the same modulus, and
that the individual polynomials \(P_i(T)\) lie in \(\mathbf Z[T]\) and
are independent of \(\ell\).  Those arithmetic conclusions have not
been used here.
\end{remark}

\section{Exact dependency graph}

\[
\begin{tikzcd}[column sep=small,row sep=small]
\substack{\text{six operations, base change,}\\\text{trace, and duality}}
  \arrow[r] & \substack{\text{rank-one and}\\\text{determinant weights}}
  \arrow[r] & \substack{\text{real-sheaf}\\\text{purity}}
  \arrow[d] \\
\substack{\text{Artin--Schreier}\\\text{orthogonality}}
  \arrow[r] & \substack{\text{Fourier inversion}\\\text{and Plancherel}}
  \arrow[r] & \substack{\text{\(L^2\)-weight}\\\text{criterion}}
  \arrow[d] \\
& \text{curve estimate}
  \arrow[r] & \substack{\text{relative affine-curve}\\\text{case}}
  \arrow[d] \\
& \substack{\text{Noether normalization,}\\\text{excision, and Leray}}
  \arrow[r] & \substack{Rf_!\text{ upper}\\\text{weights}}
  \arrow[d] \\
& \substack{\text{relative Poincare}\\\text{duality}}
  \arrow[r] & \substack{\text{smooth-proper}\\\text{purity.}}
\end{tikzcd}
\]

There is no arrow from smooth-proper purity, integral Frobenius
polynomials, hard Lefschetz, a curve Riemann hypothesis, Chebotarev, or
Picard--Lefschetz back into the graph.

\begin{formalization}
\label{form:global-weight-interface}
\textbf{Terminal interfaces.}

\begin{enumerate}
\item A sheaf-level theorem:
\[
\operatorname{WeightsLE}_\iota(\cF,w)
\Longrightarrow
\operatorname{WeightsLE}_\iota(R^if_!\cF,w+i).
\]
\item A derived theorem using the convention
\(\operatorname{WeightsLE}(\cH^j\cK,w+j)\).
\item Separate reduction lemmas for quasi-finite maps, composition,
open--closed triangles, relative Noether normalization, and affine-line
projections.
\item A Frobenius-equivariant proper-support base-change isomorphism at
every closed point.
\item A mixedness lemma based on real envelopes and reduced real fiber
\(L\)-functions.
\item A spectral-sequence interface requiring only a finite filtration
and subquotient maps, never degeneration.
\item The terminal theorem
\cref{thm:deligne-laumon-direct-image}.
\item Relative duality in the exact form
\eqref{eq:relative-pd-weight}, and the terminal purity corollary
\cref{cor:smooth-proper-purity}.
\end{enumerate}

The architecture follows the reductions in
\cite{DeligneWeilII,LaumonFourier,KiehlWeissauer,StanfordL24}.
\end{formalization}


